% !PRINTED 780
\weberpage{780}
% running head "II. Nachtrag." at the top of this page is not transcribed as body text.
%<<CONT
\[
Q(s) = \frac{1}{m(s-1)} \overset{A}{\Sigma} \chi(A) + \overset{A}{\Sigma} \chi(A)\, C(s),
\]
and from this the theorem:

2. \gesperrt{The sums $Q_2, \ldots, Q_\mu$ take finite values for $s=1$,}
$Q_1$ \gesperrt{becomes infinite, and indeed we have}
\[
\mathrm{Lim}_{s=1} (s-1)\, Q_1(s) = \frac{\mu}{m}.
\]

In order to transform the sums $Q(s)$, we denote by $p$ every prime that does
not divide $m$, and multiply all the infinite series of the form
\[
1 + \frac{\chi(p)}{p^s} + \frac{\chi(p^2)}{p^{2s}} + \cdots
= \frac{1}{1 - \chi(p)\, p^{-s}}
\]
with one another. In this way we obtain [cf. \S.\,192, (8)]:
\weq{11}{Q(s) = \overset{p}{\Pi} \frac{1}{1 - \chi(p)\, p^{-s}}.}
Now if $p$ belongs to the exponent $f$, that is, if $f$ is the least positive
number satisfying the congruence
\weq{12}{p^f \equiv 1 \pmod{m}}
and if $\mu = ef$, then among the $\chi(p)$ every $f$th root of unity occurs
exactly $e$ times, and we have therefore (for a variable $x$):
\[
\overset{\chi}{\Pi} [1 - \chi(p)\, x] = (1 - x^f)^e,
\]
and accordingly there follows from (11):
\weq{13}{Q_1 Q_2 \ldots Q_\mu = \overset{p}{\Pi} \frac{1}{(1-p^{-fs})^e}.}

We now introduce, without assuming the irreducibility of the cyclotomic
equation, the cyclotomic field $\Omega_m$, whose still unknown degree we shall
denote by $\nu$. But since the primitive $m$th root of unity $r$ satisfies a
rational equation of degree $\mu$, we have
\weq{14}{\nu \leqq \mu.}

The numbers $\omega$ of the field $\Omega_m$ can all be represented in the form
\weq{15}{a\omega = a_0 + a_1 r + \cdots + a_{\nu-1} r^{\nu-1}}
with integral rational coefficients $a, a_0, a_1, \ldots, u_{\nu-1}$; % [sic] -- printed $u_{\nu-1}$; display (15) has $a_{\nu-1}$
and if
$\omega$ is an \gesperrt{integer}, we can set for $a$ a fixed integer, a closer
knowledge of which is not neces-
%CONT>>
